\section{A common potentially crystalline family at a type prime}

% BEGIN K84P1
\paragraph{K84P1. The local comparison on the two places of a type label.}
Use the native hypotheses of $\mathrm{StepTypeAlt}$, including
$\mathrm{admDZ}(P)$, the three automorphic and Galois inputs
$\mathrm{adm1},\mathrm{adm2}_{\rm alt},\mathrm{adm5}$ and the
comparison laws $\mathrm{tie0}$. Fix $\ell_i=g_k=p$, an automorphic
representative above $x$, and an occurring type with squarefree label
$s(X)$ avoiding both designated units. The depth-zero test specifies
one of the two units $u$ and gives at the distinguished split place
the degree $n_V=r+t$ label $s(X)(X-u)^t$, with zero monodromy.
At its conjugate place it gives the multiset of inverses of these
roots, again with zero monodromy and full degree. The constant term
is nonzero by the occurring-type label condition. The calculation of
P84C3 applies to this individual label: the tame generator
distinguishes the nonzero root characters, and only $u$ (respectively
$u^{-1}$) can occur twice. Its order divides $p-1$. This gives the
noncritical condition $\mathrm{Nc}$ at both places. No goodness
condition on other points is needed.

Choose a base-change list for this representative by $\mathrm{adm1}$.
At the tested split place its local sum is the indicated parameter.
Noncriticality passes to each summand. Wild and zero-monodromy
conditions pass by direct sum, inertia polynomials multiply, and
the repetition threshold in the definition is nonincreasing with
dimension. An eigenvalue reaching its summand's threshold therefore
reaches the total threshold. The real unramified twists preserve
all these conditions. The centered cuspidal realizations of
$\mathrm{adm2}_{\rm alt}$, with the algebraic twists required
by regularity of the list as in C84F5, now give a de Rham sum. More explicitly, the top
member of each length-$d$ string has its infinity weights in the
regular total weight shifted by $(d-1)/2$. This fixes the parity
of the algebraic twist. Undo the common centering and take the
$d$ successive Tate shifts of this one realization; direct sums
and these twists preserve de Rhamness.
On split unramified Frobenius elements its polynomial agrees with
that of the good $r_x$, because both are the normalized Hecke
polynomial of the given point. Density and semisimplicity identify
the two representations over the algebraic closure of the
coefficient field, so the de Rham and Weil--Deligne comparisons
transfer by $\mathrm{tie0}$.

For clarity, inertia agreement in the weakened cuspidal comparison
has a guard. On a commutator of inertia, each nonzero-root character
of the label is one. The full product polynomial there is therefore
$(X-1)^{n_V}$, and the corresponding polynomial for each summand
is $(X-1)$ to that summand's dimension, by monic divisibility.
Unramified normalizing twists have the same inertia polynomials.
Together with zero monodromy these facts meet the commutator
guard in the comparison. In degree one the comparison is the
character identity. On the other summands, apply the guarded
inertia identity and multiply over the direct sum and its Tate
shifts. Thus at either place $w$ the Galois parameter has zero
monodromy and for every $h\in I_w$ its inertia polynomial is
\[
              \prod_{y\in M}(X-\omega_{k,w}(y,h)),
\]
where $M$ is the specified root multiset, inverted at the conjugate
place. The same chosen $u$ serves both places. This is
$\mathrm{TypeLGC}(P)$ with the required quantifiers.
% END K84P1

% BEGIN K84P2
\paragraph{K84P2. The common field and the bounded set of tame labels.}
Assume now the native premises of $\mathrm{StepPotDirect}$, in
particular $\mathrm{TypeLGC}(P)$, the local bridge and character
ties, the crystalline restriction and tensor comparisons and
$\mathrm{TieInePoly}$. Fix its $i,k,\pi,X_s,w$.
The type locus has a single tuple of squarefree labels. Since
$\pi\in X_s$, each local label occurs and hence is monic of
degree $r$, with nonzero constant term. At the distinguished
member of the pair let $c$ be the designated residue unit, and use
\[
           Y=\{\text{roots of }s\}\cup\{c,c^{-1}\}.
\]
At the conjugate place invert this set. The choice of $u$ in
K84P1 may vary with the point, but every resulting multiset is
supported in $Y$. The separate character tie supplies the label
$u_0\in Y$ of the actual de Rham line $\chi_1$ at this $w$,
including inversion at the conjugate place. Its nilpotent
monodromy is zero by dimension one. Thus each $r_x$ and the line
are potentially crystalline, by the de Rham and zero-monodromy
comparison transported through the coefficient bridge.
The full inertia polynomial on the line is
$X-\omega_{k,w}(u_0,h)$ by the one-dimensional label identity.

Write $m_G$ for the product group's order in the datum.
Every $y\in Y$ satisfies $y^{m_G}=1$: a root of degree
$d\le r$ has order dividing $p^d-1$, and
$|\mathrm{GL}_r(\mathbb F_p)|=\prod_{j=0}^{r-1}(p^r-p^j)$
contains the factor $p^{r-d}(p^d-1)$. The designated units have
$d=1$, and inversion preserves the assertion. Hence all their
orders are prime to $p$. Moreover $|Y|\le r+2\le n_V+1$.

The type prime splits completely, so the completion $K$ has
degree one over $\mathbb Q_p$. For each point take a finite
extension over which its representation becomes crystalline; do
likewise for $\chi_1$. For each $y\in Y$ the open-kernel clause of
the tame label supplies an open subgroup on whose inertia its
character is trivial. Via the decomposition-group bridge this
contains the inertia of a finite local extension: an open
neighbourhood of the identity in inertia contains the intersection
with an open subgroup of the whole local Galois group. Take the
normal closure $K'$ of the compositum of these fields and $K$.
There are finitely many points and labels. Crystallinity persists
over $K'$ by the given restriction comparison. Normality makes
$K'_0=K'\cap\mathbb Q_p^{\rm nr}$ stable, with the prescribed
action on unramified constants. The same $K'$ works on all points
and on the one-dimensional character, and it kills the inertia
characters attached to $Y$ on $I_K\cap G_{K'}$.
% END K84P2

% BEGIN K84P3
\paragraph{K84P3. Recovering simultaneous characters on crystalline periods.}
Put $H=G_{K'}$ and $I=\ker(G_K\to\mathrm{Aut}(K'_0))$.
Then $I/H$ is the image of $I_K$. Indeed the fixed field in $K'$ of the image of $I_K$ is
$K'\cap K^{\rm nr}=K'_0$, since $K=\mathbb Q_p$.
The comparison on period-ring constants identifies this $I$ with
the inertia group of the specified period datum. The inertia
identification of the bridge reads $I_K$ as the designated
global inertia.

Lift $\omega_{k,w}(y,h)$ to the coefficient DVR by the supplied
primitive $m_G$th root and then map it into the local coefficient
field $L$. This lift exists and is unique: the value is an
$m_G$th root by multiplicativity in $y$, and the powers of the
chosen primitive root have distinct complex images exhausting
those roots. The two embeddings of the DVR are the given ones.
The lifts are multiplicative in $y$ and in $h$. On the cyclic
group generated by $Y$ they kill $I_K\cap H$ by K84P2. They
therefore descend uniquely to characters $a_y:I/H\to L^\times$.
Each has order dividing the order of $y$. Moreover all are
powers of one character: take a generator of the finite cyclic
group generated by $Y$ in $\overline{\mathbb F}_p^\times$,
and apply the same multiplicative lift.

For any one of our representations let $M$ be its specified
multiset and $D'=D_{\rm cris,K'}$. The action of $I$ on $D'$
is $L$-linear and factors through the finite group $I/H$.
For each $g\in I$ choose $h\in I_K$ in its coset modulo $H$.
The polynomial with roots the lifted $\omega_{k,w}(y,h)$,
$y\in M$, has coefficients in the original DVR and its complex
image is the prescribed inertia polynomial. By the
$\mathrm{TieInePoly}$ premise the local image of this polynomial
annihilates every crystalline period under $h$. On these periods
$h$ and $g$ act identically. Thus the action of $g$ is annihilated
by $\prod_{y\in M}(X-a_y(g))$, with the multiplicities in $M$.

Let $C$ be the cyclic image of the single character generating
the $a_y$. An element in its kernel has unipotent action by the
last polynomial. Since it has finite order on the period space
in characteristic zero, its action is the identity. On $D'$
the action now factors through $C$. Choose a preimage of a
generator. Its minimal polynomial divides both the above
product and the separable polynomial $X^{|C|}-1$; hence it is
diagonalisable over $L$ with eigenvalues among the listed
$a_y$ evaluated at that generator. Every eigenspace is a
simultaneous character space for the entire cyclic quotient.
This proves the full period eigenspace sum, the triviality on
$H$, and the prime-to-$p$ orders required by
$\mathrm{InertiaThrough}$. For $\chi_1$ the allowed set is
the singleton $\{a_{u_0}\}$ throughout.
% END K84P3

% BEGIN K84P4
\paragraph{K84P4. A Frobenius-stable support for the twisted family.}
By the crystalline tensor and dual comparison over the common
$K'$, $r_x\chi_1^{-1}$ is crystalline there, and every character
of its period inertia action belongs to
$S=\{a_ya_{u_0}^{-1}:y\in Y\}$. The evaluation pairing of the
line with its inverse divides the inertia eigenvalue by
$a_{u_0}(g)$ for every $g$. Thus $1\in S$ and $|S|\le n_V+1$.

Take a lift $F\in G_K$ of the required residue-field Frobenius on
$K'_0$. It exists because $K=\mathbb Q_p$ and the automorphism of
this unramified field extends to the algebraic closure. The action
of $F$ on $D_{\rm cris,K'}$ preserves periods by normality
and is $L$-linear, with its specified semilinear action on $K'_0$.
For a nonzero vector in the $a$-space, its image under $F$ belongs
to the character $g\mapsto a(F^{-1}gF)$, since $I$ is normal.
It is nonzero and therefore this character occurs in the same
representation and belongs to $S$. Take $\Sigma$ to be $1$
together with the characters that actually occur at any retained
point. The preceding argument makes $\Sigma$ Frobenius-stable,
and K84P3 shows that it gives every required inertia decomposition.
As $\Sigma\subset S$, stability has incurred no cardinality loss.
These $K',F,\Sigma$ give $\mathrm{PotFam}$ for the whole direct
family. This proves $\mathrm{StepPotDirect}$ from its stated
inputs. T84L5 supplies the Tate-dual family with the same field,
and provides the two signed jump tests under the host hypothesis.
The quantitative lifting theorem $\mathrm{admT}(C_T)$ enters
later in T84L1.
% END K84P4
